\documentclass[12pt]{article}
\usepackage[utf8]{inputenc}
\usepackage[T1]{fontenc}
\usepackage[margin=1in]{geometry}
\usepackage{setspace}
\usepackage{xcolor}
\usepackage[colorlinks=true, urlcolor=blue!60!black, linkcolor=blue!60!black]{hyperref}
\setlength{\parindent}{2em}
\setlength{\parskip}{1.25em}
\renewcommand{\baselinestretch}{1.0}
\begin{document}
\begin{center}
{\LARGE \textit{South-Central Timber Development, Inc.\ v.\ Wunnicke}, 467 U.S.\ 82 (1984)\par}
\vspace{0.8em}
{\large Case Brief --- Four-Page Edition\par}
\end{center}
\vspace{1em}

\section{Memory Jogger}

Alaska sold state timber on condition that buyers process it in Alaska before export. A parallel federal-lands policy was not ``unmistakably clear'' congressional authorization, and a plurality held that the market-participant exception does not let a State impose conditions regulating a downstream market in which it does not participate (\href{https://supreme.justia.com/cases/federal/us/467/82/}{Justia}).

\section{Detailed Case Facts}

In 1980 Alaska announced a sale of about 49 million board-feet of state-owned timber near Icy Cape. Under state regulation and a 1974 Governor's Policy Statement, the contract required ``primary manufacture'' in Alaska before the timber could be shipped out of state. The condition applied only to state timber and aimed to protect existing industries, foster new ones, raise revenue, and manage forests. Because it depressed bids, Alaska charged a lower price. Buyers usually complied by sawing logs into ``cants,'' which wasted wood and were worth less than round logs. South-Central Timber, an Alaska logging company without a sawmill, exported unprocessed logs almost entirely to Japan, which took over 90\% of Alaska timber. Since 1928, federal rules had restricted export of unprocessed timber from \emph{federal} lands in Alaska.

\section{Procedural History}

The district court enjoined the clause under the dormant Commerce Clause. The Ninth Circuit reversed, holding that federal primary-manufacture policy for federal lands ``implicitly authorized'' Alaska's parallel policy. The Supreme Court granted certiorari, with the United States supporting South-Central, and reversed on May 22, 1984.

\section{Judicial Votes}

Reversed and remanded, \textbf{6--2} (Marshall not participating).
\begin{itemize}
\item \textbf{Opinion of the Court (Parts I--II):} Justice White, joined by Chief Justice Burger and Justices Brennan, Blackmun, Powell, and Stevens.
\item \textbf{Plurality (Parts III--IV):} Justice White, joined by Justices Brennan, Blackmun, and Stevens.
\item \textbf{Concurrences:} Justice Brennan; Justice Powell, joined by Burger (in part and in the judgment).
\item \textbf{Dissent:} Justice Rehnquist, joined by Justice O'Connor.
\end{itemize}

\section{Holding}

\textbf{Majority:} To shield a state restraint from the dormant Commerce Clause, congressional intent must be ``unmistakably clear''; a federal policy for federal lands does not authorize a comparable restraint on state lands. \textbf{Plurality:} The market-participant doctrine lets a State burden commerce only in the market in which it participates, not through conditions with a substantial regulatory effect downstream (here, timber processing). Alaska's local-processing mandate is virtually \textit{per se} invalid, reinforced by heightened Foreign Commerce Clause scrutiny.

\section{Analysis of Opinions}

\textbf{Majority (White), Part II.} \textit{Sporhase} and \textit{New England Power} require that Congress affirmatively contemplate the state restraint. ``Expressly stated'' is not talismanic, but intent must be unmistakably clear, ensuring that any authorization of Balkanizing barriers is a national decision, especially where foreign relations are affected. Congress legislated only for federal lands.

\textbf{Plurality, Parts III--IV.} \textit{Alexandria Scrap} left sellers free to forgo the bounty; \textit{Reeves} lacked foreign commerce, a natural resource, and post-sale restrictions; and in \textit{White v.\ Massachusetts Council} everyone was ``working for the city,'' whereas Alaska controlled buyers' dealings with a ``stranger to the contract'' after payment. Drawing on restraints on alienation and antitrust limits on vertical restraints, the plurality rejected Alaska's argument that it could reach the same end by other means. A naked local-processing rule is virtually \textit{per se} invalid under \textit{Pike}, \textit{Foster-Fountain}, and \textit{City of Philadelphia}, and the Nation must ``speak with one voice'' in foreign trade.

\textbf{Brennan (concurring).} The dispute exposed ``the inherent weakness'' of the market-participant doctrine.

\textbf{Powell (concurring in part).} He would have remanded the remaining issues.

\textbf{Rehnquist (dissenting).} The participant-regulator line is ``both artificial and unconvincing.'' Alaska was paying buyers, through a lower price, to hire local processors, functionally identical to permissible alternatives.

\textbf{Later.} \textit{Maine v.\ Taylor} (1986) applied the clear-statement canon, and courts treat the downstream limit as the doctrine's outer boundary.

\section{Examples: Future Applications}

\begin{enumerate}
\item \textbf{Same-Side: State Lithium Refining.} Nevada requires state-land lithium lessees to refine in-state, citing a federal-lands rule. That is not clear authorization, and it regulates the downstream refining market.
\item \textbf{Same-Side: Cement Resale Limits.} A state cement plant bars distributors from reselling outside the State after title passes, a post-sale restriction of the kind absent in \textit{Reeves}.
\item \textbf{Opposite-Side: A State Refinery's Customers.} A state-owned refinery selling its output only to local battery makers is choosing customers in its own market, as in \textit{Reeves}.
\item \textbf{Opposite-Side: Cash Subsidies.} Paying in-state refiners from the general fund, while lessees ship anywhere, leaves buyers their choice.
\item \textbf{Fence-Sitter: Resident Logging Crews.} Requiring Alaska crews to cut and haul timber off state land may be inside the timber-sales market or downstream of it.
\end{enumerate}

\section{Critique}

\textbf{Market definition.} The plurality defined the market by ``intuition,'' without an objective test. \textbf{Formalism versus transparency.} Rehnquist argued that Alaska could reach the same result by integration or subsidies; yet the line forces visible appropriations rather than hidden auction discounts. \textbf{Fragility.} Only four Justices joined Parts III--IV.

\section{Key Quotations}

\begin{quote}
``There is no talismanic significance to the phrase `expressly stated,' however; it merely states one way of meeting the requirement that, for a state regulation to be removed from the reach of the dormant Commerce Clause, congressional intent must be unmistakably clear. The requirement that Congress affirmatively contemplate otherwise invalid state legislation is mandated by the policies underlying dormant Commerce Clause doctrine.''
\end{quote}
\noindent --- Justice White, for the Court, 467 U.S.\ at 91--92.

\begin{quote}
``The limit of the market participant doctrine must be that it allows a State to impose burdens on commerce within the market in which it is a participant, but allows it to go no further. The State may not impose conditions, whether by statute, regulation, or contract, that have a substantial regulatory effect outside of that particular market. Unless the `market' is relatively narrowly defined, the doctrine has the potential of swallowing up the rule that States may not impose substantial burdens on interstate commerce even if they act with the permissible state purpose of fostering local industry.''
\end{quote}
\noindent --- Justice White, plurality opinion, 467 U.S.\ at 97--98.

\begin{quote}\sloppy
``The fact that the state policy in this case appears to be consistent with federal policy---or even that state policy furthers the goals we might believe that Congress had in mind---is an insufficient indicium of congressional intent. Congress acted only with respect to federal lands; we cannot infer from that fact that it intended to authorize a similar policy with respect to state lands.''
\end{quote}
\noindent --- Justice White, for the Court, 467 U.S.\ at 92--93.

\end{document}
